\section{Discussion}
\label{sec:discussion}

\subsection{Methodological Contributions}

Our approach is complementary to ongoing theoretical work on multiparameter persistence \citep{BotnanLesnick2022}, which addresses scale conflation algebraically but faces computational barriers for large point clouds. We take a practical, domain-informed route: hierarchical clustering at validated cutoff distances decomposes the point cloud by organisational level before computing standard single-parameter persistent homology at each level.

The cutoff distance, which we term \emph{domain-informed} rather than ``goal-dependent'' (to avoid ambiguity with the sport-specific meaning of ``goal''), defines the organisational level under analysis. The three validated regimes (individual, tactical, team) emerge from systematic parameter sweep and reflect genuine hierarchical structure in the competitive system. The tactical cutoff (12.0~m) is explicitly a domain-informed choice, justified by (i) its position within the automated metric range (6.87--16.31~m), (ii) its correspondence to standard football zone dimensions, and (iii) sensitivity analysis showing robust \(H_1\) detection across the plausible range (Section~\ref{sec:sensitivity}).

The adaptive filtration formula addresses a practical coupling between the clustering and filtration steps. The P75 threshold is empirically justified by ablation showing it is the minimum percentile achieving full detection sensitivity (Section~\ref{sec:sensitivity}). This is a general solution applicable to any multi-scale TDA pipeline that uses clustering as a preprocessing step.

\subsection{Multi-Scale Topological Structure}
\label{sec:scale_complementarity_discussion}

The complementarity of the individual and tactical scales is supported by quantitative evidence (Section~\ref{sec:scale_complementarity}): \textbf{Spearman} rank correlation between individual-scale and tactical-scale \(H_1\) (count-based) \(\rho=0.254\), \(p<0.001\); 95\% match-level bootstrap CI \([0.200,\,0.314]\) (1{,}000 resamples over matches). \textbf{Fisher's exact test} on the \(2\times2\) presence/absence table gives \(\mathrm{OR}=3.52\), \(p=0.093\), with 95\% match-level bootstrap CI \([2.59,\,13.53]\). Spearman \(\rho\) remains the primary summary for scale complementarity. The weak positive correlation likely reflects the confound that dense player configurations create opportunities for topological features at multiple scales simultaneously. Despite this, the scales capture predominantly distinct structural information, as demonstrated by divergent multi-match frame \(H_1\) presence rates ($97.0\%$ vs $19.3\%$) and persistence profiles.

\paragraph{Baseline versus topology --- non-redundancy without held-out predictive utility.}
Two complementary statements characterise the relationship between tactical \(H_1\) total persistence and the four geometric baselines. \emph{Non-redundancy (correlational; Section~\ref{sec:baseline}).} Spearman \(\rho\) is \(-0.005\), \(0.356\), \(0.431\), and \(0.126\) for length, width, hull area, and Voronoi entropy (\(p=0.857\), \(p<10^{-40}\), \(p<10^{-65}\), \(p=9.9\times10^{-7}\)), and the partial \(R^2\) of tactical-scale \(H_1\) after the other three baselines is \(0.025\), \(0.036\), \(0.091\), \(0.004\); these values, in particular the \(0.091\) for hull area, demonstrate that tactical \(H_1\) explains residual variance in the geometric block that the geometric features themselves do not account for. \emph{Predictive utility (cross-validated; Section~\ref{sec:predictive}).} The natural follow-up question --- whether that residual variance translates into incremental predictive performance on held-out matches --- is answered negatively at the present sample size: a GroupKFold-over-matches logistic regression for the binary build-up-phase label yields AUC \(=0.693\) for the baseline block and AUC \(=0.687\) for the baseline-plus-tactical-\(H_1\) block, \(\Delta\text{AUC}=-0.005\) with match-level bootstrap 95\% CI \([-0.015,\,+0.001]\) and stratified permutation \(p=1.0\). The same null finding holds when tactical-\(H_1\) total persistence is replaced by \(H_1\) count or max-persistence and when the target is replaced by the chaotic-phase indicator (Table~\ref{tab:predictive}). The two statements are not in contradiction. Non-redundancy is a variance-decomposition property of the within-sample regression, whereas predictive utility on held-out matches additionally requires that the residual variance generalises between matches with only 10 training units to estimate the topology's regression coefficient. With the baseline block already achieving an AUC near \(0.7\) on this task, the marginal slot available for tactical \(H_1\) is small, and the present 10-match cohort does not have the statistical power to fill it.

The honest reading is therefore that tactical \(H_1\) carries non-redundant geometric information at the per-frame level on this sample (Section~\ref{sec:baseline}) but that the present 10-match cohort is not large enough to convert that into a detectable held-out predictive gain for phase classification (Section~\ref{sec:predictive}); the fully-specified ANCOVA with topology as the tested block at full-season scale is the definitive test of incremental predictive utility, and the present null is the effect-size benchmark a population-scale replication needs to beat.

\subsection{Temporal Dynamics and Event Correlation}

The temporal evolution result is more nuanced than initially reported: the primary match shows a large but non-significant tactical-scale persistence increase (+78.8\%, \(p=0.066\)), but the direction is not consistent across matches. The non-significant result for the primary match, despite a large effect size, reflects the low statistical power inherent in single-match half-level comparisons ($n=10$ and 9 windows with non-zero tactical persistence per half). Half-level persistence dynamics appear driven by match-specific factors such as tactical adjustments, scoreline effects, and substitutions, rather than universal trends.

\paragraph{Multi-match formal test (Section~\ref{sec:temporal}).}
A linear mixed model on per-window \emph{tactical-scale} $H_1$ persistence over all 10 matches ($n=354$ two-minute windows) with match as a grouping factor and ``half'' as a fixed and random-slope effect gives: fixed effect $\hat\beta_1$ (second half) $=-0.081$, 95\% CI $[-0.172,\,0.0093]$, LMM $p=0.079$; estimated random-slope variance $\widehat{\mathrm{Var}}(u_1)=0.00101$; stratified permutation $p=0.051$ (10{,}000 within-match permutations of the half label). These estimates align the informal ``match-specific dynamics'' phrasing with evidence that is borderline at $\alpha=0.05$, rather than a uniform half effect across matches.

The real event correlation, based on 104{,}722 event--topology pairs across 10 matches, provides the first evidence that persistent homology features are genuinely responsive to match dynamics. The coherent pattern is that disruption decreases persistence whilst organisation increases it; this validates the topological interpretation and suggests applications in real-time tactical monitoring. Event-window sensitivity (Table~\ref{tab:event_window}) shows the same sign for the five headline tactical-scale event types from $\pm0.5$~s to $\pm5$~s.

\subsection{Bilateral Coupling}
\label{sec:bilateral_discussion}

The bilateral decomposition (Section~\ref{sec:bilateral}, Tables~\ref{tab:bilateral_perteam}--\ref{tab:bilateral_coupling}) makes the team identity that the 22-player analysis discards an explicit object of inference. Two findings are worth highlighting. First, per-team tactical \(H_1\) presence rate (\(\approx 36\%\) for each of home and away) is roughly twice the 22-player merged rate (\(19.3\%\pm 7.2\%\); Table~\ref{tab:h1multi}); the merge therefore obscures formation structure that is recovered cleanly by splitting on team identity, and the per-team rate is symmetric between home and away to within 0.3~percentage points (a sanity check that the framework treats the two teams identically). Second, the per-frame cross-team coupling is small in every metric we examined: Spearman \(\rho=0.037\) at lag 0 with 95\% match-level bootstrap CI \([-0.018,\,0.087]\) (i.e.\ not distinguishable from zero); odds ratio for joint \(H_1\) presence \(1.32\) (only marginally above the independence value of \(1\)); and lagged \(\rho\) staying below \(0.06\) in absolute value out to \(\pm 10\) frames. The two teams' tactical persistence sequences therefore evolve close to independently at frame resolution within these 1{,}500 frames. This is itself a non-trivial finding --- it sets the per-frame bilateral coupling at zero as a null reference, against which any later evidence for a tactical-context-dependent coupling (during pressing sequences, during sustained possession, around set plays) can be tested. The bilateral coupling statistic is the first quantitative anchor for the bilateral topological-coupling hypothesis that motivates the population-scale extensions of \citet{Brown2026inprep}; the present null result is precisely the kind of effect-size benchmark the full-season analysis needs to power.

\subsection{Limitations}
\label{sec:limitations}

\textbf{Linkage selection is a substantive methodological choice.} The chaining effect at the tactical scale (Section~\ref{sec:linkage}) yields 153 tactical \(H_1\) loops under single-linkage versus 923 (complete-linkage) and 936 (Ward's method) over the same 600-frame, 4-match comparison sample. This is nearly an order of magnitude and is large enough that the linkage choice changes the scientific picture rather than merely the precision of a conservative estimate; we should therefore treat it as a methodological decision, not a conservatism preference. We retain single-linkage in this paper for the domain reason given in Section~\ref{sec:linkage} (clusters as chain-connected proximity groups, robust against partition artefacts), but we regard the order-of-magnitude gap as an open scientific question rather than a settled one.

A principled criterion for linkage selection should come from the dynamical system rather than from clustering quality metrics alone. A natural candidate is to choose the linkage method under which the discovered governing equation for a persistence functional --- for example the landscape \(L^2\) norm \(t \mapsto \|\lambda_\delta(t)\|_{L^2}\) or the tactical \(H_1\) total-persistence sequence --- is most sparse, in the sense of sparse identification of nonlinear dynamics (SINDy; \citealp{Brunton2016}). Under this view, the correct linkage is the one for which the post-clustering point cloud is the most natural coordinate system for the underlying dynamics, with sparsity of the recovered library acting as an information-theoretic selection criterion. Operationalising this criterion (i.e.\ choosing the dynamical state, defending sparsity-of-library against AIC/BIC/cross-validation baselines, and confirming that the chosen linkage's persistence sequence is genuinely lower-dimensional) is beyond the scope of the present validation paper and is deferred to the forthcoming full-season work \citep{Brown2026inprep}. The 600-frame linkage table in Section~\ref{sec:linkage} should be read as setting up this question, not resolving it.

\textbf{Tactical cutoff is a domain-informed choice within an automated-metric range.} The silhouette-optimal cutoff (16.31~m) and the information-content-optimal cutoff (6.87~m) bracket the chosen 12.0~m. The half-zone-width justification (Section~\ref{sec:cutoff}) is domain-reasonable but ultimately subjective, and Section~\ref{sec:sensitivity} reports the operative range (\(\delta \in [6, 14]\)~m) on the 150-frame primary-match sample as an empirical sensitivity rather than a derivation of \(\delta\). A more principled criterion would replace the domain heuristic with a stability property of the dynamical persistence sequence: the most stable tactical cutoff is the one minimising the variation of the persistence landscape path under small perturbations of \(\delta\),
\begin{equation}
  \delta^{\star} \;=\; \arg\min_{\delta}\;
    \mathrm{Var}_{\epsilon}\bigl\|\,\lambda_{\delta+\epsilon}(\cdot) - \lambda_{\delta}(\cdot)\,\bigr\|_{L^{2}},
  \label{eq:landscape_stability_criterion}
\end{equation}
where \(\epsilon\) ranges over a small neighbourhood of zero and \(\lambda_\delta(\cdot)\) is the per-match persistence landscape path \(t\mapsto\lambda_\delta(t)\) \citep{Bubenik2015}. This is a computable quantity given the GUDHI landscape infrastructure already used in Section~\ref{sec:scale_complementarity} (Table~\ref{tab:tda_distances}; \citealp{Tauzin2021}) and is the natural successor to the present domain-informed choice. Evaluating it on a single match would be premature --- the sample size for landscape variance estimation is the number of matches, not the number of frames --- so we defer evaluation to the persistence-landscape companion paper flagged in Section~\ref{sec:outlook}, treating the 12.0~m value reported here as a domain-informed reference within the empirically-validated \([6, 14]\)~m range.

\textbf{Two-scale \(H_1\), not three.} Our multi-scale framework identifies three \(H_0\) regimes but only two \(H_1\) regimes. Team-scale \(H_1\) absence is a structural consequence of the 1--2 centroids produced at \(\delta=30.0\)~m, not a methodological limitation. Effective team-scale \(H_1\) analysis would require alternative representations (e.g.\ spatial density fields or Delaunay triangulations) rather than centroid-based point clouds.

\textbf{Broadcast tracking resolution.} All data in this study are broadcast-derived (SkillCorner, 10~Hz). Surveys of tracking in team sports highlight trade-offs among sampling rate, spatial precision, and deployment cost \citep{Cummins2013, Aughey2011}; broadcast products typically offer lower spatial precision than higher-frequency optical systems (e.g.\ 25~Hz), which may affect persistence magnitude estimates. The framework's structural findings (scale regimes, $H_1$ presence rates, sensitivity profiles) are expected to be robust across tracking technologies, but direct persistence magnitude comparison with optical data remains for future work.

\subsection{Outlook}
\label{sec:outlook}

Two extensions of this work are under active development, and each directly resolves one of the methodological limitations of Section~\ref{sec:limitations}. First, the 10-match evidence base presented here is being scaled to a full Championship season ($\sim$540 matches) to characterise population-level distributions of $H_0$/$H_1$ counts, barcode lengths, and landscape norms, and to test tactical-fingerprint classification at population scale \citep{Brown2026inprep}; this is the natural setting for the SINDy-on-persistence linkage criterion of Section~\ref{sec:limitations}, because sparse identification of a governing equation for $t\mapsto\|\lambda_\delta(t)\|_{L^2}$ requires the population-sized sample of sequences that a full season provides. Second, persistence landscape dynamics \citep{Bubenik2015,Chazal2014} are being developed to treat each match as a path $t \mapsto \lambda_\delta(t)$ in landscape space, enabling stability-regime characterisation and change-point detection within and across matches; this is the natural setting in which the landscape-stability cutoff criterion of equation~\eqref{eq:landscape_stability_criterion} can be evaluated, since the resampling unit for $\mathrm{Var}_\epsilon\|\lambda_{\delta+\epsilon}-\lambda_\delta\|_{L^2}$ is the match-level landscape rather than the per-frame summary. These extensions address the main open questions emerging from the present results: whether tactical-scale $H_1$ presence is a population-stable feature, whether SINDy-recovered governing equations select between candidate linkage methods, and whether functional representations of persistence dynamics recover tactical structure beyond what static diagrams capture.
